\subsection{凸函数}

定义 5. --- 设 X 是仿射空间 E 的一个凸子集。定义在 X 上的实值有限函数称为凸的（相应地严格凸的），如果对 X 中任意两个不同的点 x、y 及任意实数 \(\lambda\)（\(0 < \lambda < 1\)），我们有：

\[
f (\lambda x + (1 - \lambda) y) \leqslant \lambda f (x) + (1 - \lambda) f (y)\tag{1}
\]

（相应地

\[
f (\lambda x + (1 - \lambda) y) <   \lambda f (x) + (1 - \lambda) f (y)).\tag{2}
\]

当 E = R 时，凸函数的这一定义与 FVR, I, 第 32 页 中的定义相同。进而，f 在 X 中是凸的（相应地严格凸的），当且仅当对每条仿射直线 \(D \subset E\)，f 在 \(X \cap D\) 上的限制在 \(X \cap D\) 中是凸的（相应地严格凸的）。

例. --- 若 \(f\) 是 \(\mathbf{E}\) 上的一个仿射线性函数，则 \(f\) 与 \(f^2\) 都是 \(\mathbf{E}\) 上的凸函数；对 \(f\) 而言这是显然的，因为

\[
f (\lambda x + (1 - \lambda) y) = \lambda f (x) + (1 - \lambda) f (y);
\]

另一方面，若 \(\alpha = f(x)\)、\(\beta = f(y)\)，则；

\[
\lambda \alpha^ {2} + (1 - \lambda) \beta^ {2} - (\lambda \alpha + (1 - \lambda) \beta) ^ {2} = \lambda (1 - \lambda) (\alpha - \beta) ^ {2} \geqslant 0
\]

其中 \(0 < \lambda < 1\)；进而，\(f^2\) 在仿射直线 \(\mathbf{D} \subset \mathbf{E}\) 上的限制是严格凸的（若 \(f|\mathbf{D}\) 不是常数）。

实值函数 \(f\)（定义在 \(\mathbf{X}\) 上的）称为凹的（相应地严格凹的），如果 \(-f\) 是凸的（相应地严格凸的）。也就是说，对任意两个不同的点 \(x, y\)（\(\mathbf{X}\) 中的）及每个数 \(\lambda\)（满足 \(0 < \lambda < 1\)），我们有

\[
f (\lambda x + (1 - \lambda) y) \geqslant \lambda f (x) + (1 - \lambda) f (y)
\]

（相应地

\[
f (\lambda x + (1 - \lambda) y) > \lambda f (x) + (1 - \lambda) f (y)).
\]

X 到 R 中的映射称为仿射的，如果它既是凸的又是凹的（cf.~II, 第 78 页, exerc. 11）。

命题 19. --- 设 X 是仿射空间 E 中的凸集；f 是定义在 X 上的实值函数。用 F（相应地 F'）表示点 \((x, a) \in \mathbf{E} \times \mathbf{R}\)（满足 \(x \in X\) 且 \(f(x) \leqslant a\)（相应地 \(x \in X\) 且 \(f(x) < a\)））组成的集。则下列条件等价：

\begin{enumerate}
\def\labelenumi{\alph{enumi})}
\item
  函数 \(f\) 是凸的。
\item
  集 \(\mathbf{F}\) 在仿射空间 \(\mathbf{E} \times \mathbf{R}\) 中是凸的。
\item
  集 \(\mathbf{F}'\) 在仿射空间 \(\mathbf{E} \times \mathbf{R}\) 中是凸的。
\end{enumerate}

我们证明 \(a) \Rightarrow c)\)。设 \((x, a)\) 与 \((y, b)\) 是 \(F'\) 的两点且 \(0 < \lambda < 1\)，则 \(f(x) < a\)，\(f(y) < b\)，而若 \(f\) 是凸的，则

\[
f (\lambda x + (1 - \lambda) y) \leqslant \lambda f (x) + (1 - \lambda) f (y) <   \lambda a + (1 - \lambda) b
\]

这表明点 \(\lambda(x, a) + (1 - \lambda)(y, b)\)（\(\mathbf{E} \times \mathbf{R}\) 中的）属于 \(\mathbf{F}'\)。于是 \(\mathbf{F}'\) 是凸的。

其次我们证明 \(c) \Rightarrow b)\)。若 \((x, a), (y, b)\) 是 \(F\) 的两点，则对每个 \(\varepsilon > 0\)，\((x, a + \varepsilon)\) 与 \((y, b + \varepsilon)\) 属于 \(F'\)，且当 \(0 < \lambda < 1\) 时，\((\lambda x + (1 - \lambda)y, \lambda a + (1 - \lambda)b + \varepsilon)\) 亦然；由 \(F\) 的定义，这就蕴含 \((\lambda x + (1 - \lambda)y, \lambda a + (1 - \lambda)b)\) 属于 \(F\)。

最后 \(b) \Rightarrow a)\)，因为（采用上面的记号）若 \((\lambda x + (1 - \lambda)y, \lambda a + (1 - \lambda)b)\) 属于 F，则

\[
f (\lambda x + (1 - \lambda) y) \leqslant \lambda a + (1 - \lambda) b
\]

只要 \(a \geqslant f(x)\) 且 \(b \geqslant f(y)\)；由此得出 (1)，从而 \(f\) 是凸的。

推论. --- 若 \(f\) 在 \(\mathbf{X}\) 中是凸的，则对所有 \(\alpha \in \mathbf{R}\)，诸 \(x \in \mathbf{X}\)（满足 \(f(x) \leqslant \alpha\)（相应地 \(f(x) < \alpha\)））组成的集是凸的。

事实上，它是 F（相应地 F'）与超平面 \(E \times \{\alpha\}\)（\(E \times R\) 中的）之交在 E 上的投影。

命题 20. --- 设 f 是定义在仿射空间 E 的凸集 X 上的凸函数。则对任意族 \((x_{i})_{1\leqslant i\leqslant p}\)（X 的 \(p\) 个点的）及任意族 \((\lambda_{i})_{1\leqslant i\leqslant p}\)（\(p\) 个实数的，均 \(\geqslant 0\)，满足 \(\sum_{i = 1}^{p}\lambda_{i} = 1\)），我们有：

\[
f \left(\sum_ {i = 1} ^ {p} \lambda_ {i} x _ {i}\right) \leqslant \sum_ {i = 1} ^ {p} \lambda_ {i} f \left(x _ {i}\right).\tag{3}
\]

若 \(f\) 是严格凸的且 \(\lambda_i > 0\)（对所有 \(i\)），则

\[
f \left(\sum_ {i = 1} ^ {p} \lambda_ {i} x _ {i}\right) <   \sum_ {i = 1} ^ {p} \lambda_ {i} f \left(x _ {i}\right),\tag{4}
\]

除非诸 \(x_{i}\) 全相等。

不等式 (3) 由上面的 II, prop. 19 及 II, 第 8 页, prop. 1 得出。设诸 \(x_{i}\) 不全相等（这蕴含 \(p \geqslant 2\)）且诸 \(\lambda_{i}\) 全 \(>0\)；则点 \(z = \sum_{i=1}^{p} \lambda_{i} x_{i}\) 至少与某个 \(x_{i}\) 不同。不妨设 \(z \neq x_{1}\)，写 \(z = \lambda_{1} x_{1} + (1 - \lambda_{1}) y_{1}\)，其中 \(y_{1} = \sum_{i=2}^{p} \frac{\lambda_{i}}{1 - \lambda_{1}} x_{i}\)。则 \(y_{1} \neq x_{1}\)，且由于 \(0 < \lambda_{1} < 1\)，由假设我们有

\[
f (z) <   \lambda_ {1} f (x _ {1}) + (1 - \lambda_ {1}) f (y _ {1}).
\]

但由 (3) 有 \(f(y_{1}) \leqslant \sum_{i=2}^{p} \frac{\lambda_{i}}{1 - \lambda_{1}} f(x_{i})\)，由此即得不等式 (4)。

